% =====================================================================
\section{v7.1: Collective StrongEngine — Architectural Symmetry Completion}
\label{sec:collective-strong-engine-v71}

\subsection{The structural defect discovered at v7.0}

v7.0's introduction of \texttt{NRMO\_Collective} resolved the empirical
Faith\_Communal gap in catastrophic worlds. However, on careful
inspection of the v7.0 architecture, a fundamental asymmetry was
revealed:

\begin{center}
\begin{tabular}{lll}
\toprule
Domain & Veto layer & Strong Engine (search) \\
\midrule
Individual governance & NRMO Core (4 or 7 rules) & $\Omega$ Full (mutation, rollout, scoring) \\
Collective governance (v7.0) & P--U rules (fixed) & \textbf{absent} \\
\bottomrule
\end{tabular}
\end{center}

The v7.0 mechanisms P--U operate as a fixed-rule layer: pooling rates,
coverage levels, target selection, all parametrically frozen. There
is no active exploration of \emph{which} collective configuration is
optimal for a given world-state. This violates the NRMO architectural
principle that veto and search should be present at every scale.

\subsection{The principle: fractal NRMO}

NRMO's foundational claim, stated since v4 and v25 reformulations, is
that the same governance--execution separation applies at every
relevant scale:

\begin{quote}
``Individual, family, lineage, nation, civilisation, world — at every
scale, NRMO defines admissibility and the Strong Engine searches
within. The principle is scale-invariant.''
\end{quote}

v7.0 implemented this for the collective veto side (P--U as
admissibility constraints on collective action) but not the search
side. v7.1 completes the symmetry by introducing
\texttt{CollectiveStrongEngine}: a search engine that operates on
collective-governance configurations under collective veto.

\subsection{Architectural placement}

For each generation, per civilisation, the full pipeline becomes:

\begin{enumerate}
\item NRMO\_Core (individual veto) constructs
  $A_t^{\text{individual}}$.
\item $\Omega$ Full (individual Strong Engine) searches
  $A_t^{\text{individual}}$, returns $a^\star_{\text{indiv}}$.
\item Project $a^\star_{\text{indiv}}$ to agents
  (Section~\ref{sec:vnext-plus-v63}).
\item \emph{New:} Collective\_Core (collective veto) constructs
  $A_t^{\text{collective}}$ (the set of admissible collective
  configurations).
\item \emph{New:} CollectiveStrongEngine searches
  $A_t^{\text{collective}}$, returns $c^\star$ (collective
  configuration).
\item Apply $c^\star$ to P--U mechanisms (set pool rates, trigger
  thresholds, target priorities, structural reallocations).
\item P--U mechanisms execute their veto/rescue cascade using
  $c^\star$ as parameters, instead of the fixed values of v7.0.
\end{enumerate}

The collective layer thus mirrors the individual layer with
identical structure: veto then search, never reversed.

\subsection{Six new mechanisms (W--AB)}

\paragraph{W. Predictive trigger.}
The Collective Engine forecasts next-generation shock probability
from the current world parameters
(\texttt{shock\_probability}, \texttt{tail\_probability},
\texttt{environmental\_drag}, current $X$, current trajectory).
The rescue budget is pre-emptively augmented when the forecast
exceeds threshold. This converts the rescue cascade from
\emph{reactive} to \emph{anticipatory}: high-risk gens receive
expanded coverage \emph{before} the failure event, not after.

\paragraph{X. Target selection optimisation (triage).}
The v7.0 P mechanism rescued lineages by uniform random sampling.
v7.1 implements explicit triage: each at-risk lineage receives a
\emph{rescue priority score} composed of (i)~knowledge contribution
(per-lineage mean education), (ii)~productivity contribution
(per-lineage mean assets), (iii)~optionality contribution
(per-lineage strategic diversity contribution to civ-wide Shannon
entropy of strategies), (iv)~vulnerability complement (rescue those
with lower self-recovery probability). The Collective Engine learns
the weighting of these four factors via MC rollout: rescue
configurations are evaluated by their effect on civ continuation
$+$ diversity, and the best weighting emerges.

\paragraph{Y. Pool reallocation.}
Each generation the Collective Engine considers reallocating surplus
between tiers: family $\to$ lineage, lineage $\to$ civ, civ
$\to$ inter-civ. Reallocations are veto-checked (cannot drain a tier
below safety floor) and selected by rollout: ``what happens to
continuation rate if I move 20\% of family pool surplus to civ pool
this gen?''. This makes the multi-tier structure responsive to
real-time pressure.

\paragraph{Z. Inter-civilisation insurance.}
Civilisations with low rivalry can enter contractual mutual
insurance: each contributes a fraction to a shared pool, drawn upon
when one party faces catastrophic shock. The Collective Engine
explores which contracts to enter, the contribution rate, and the
trigger conditions. Mutual contracts increase total insurable
capital at the cost of cross-civ dependency.

\paragraph{AA. Insurance candidate exploration.}
By mechanism analogue to $\Omega$ Full's mutation/synthesis/invention
on individual actions, the Collective Engine generates candidate
\emph{collective configurations}: new tier structures (3-tier $\to$
4-tier with an individual-level reserve, $\to$ 5-tier with
inter-civ), modified trigger conditions (event-based vs.
slope-based), alternative target priority rules. These candidates
are scored by rollout.

\paragraph{AB. Collective drift control.}
By mechanism analogue to v6.4's individual cumulative\_drift, the
collective engine tracks ``rescue dependency'': accumulated rescues
per generation, normalised by civilisation size. If rescue
dependency grows unbounded, the civilisation is drifting toward an
absorbing state where its individual lineages have ceased
self-recovery and depend entirely on the pool. The Collective Veto
constrains rescue rates to prevent this drift.

\subsection{Collective Veto (Collective\_Core)}

By NRMO principle, the collective layer also has a veto component
distinct from the search component. The four veto rules:

\begin{enumerate}
\item \textbf{Pool depletion veto.}\ No collective configuration is
admissible if it depletes any pool below the safety floor in
projection (e.g.\,family pool below $0.10 \times$ historical mean).
\item \textbf{Rescue rate ceiling.}\ No configuration is admissible
if the projected per-gen rescue rate exceeds a ceiling
($0.30$ of population). Above this rate, individual incentive
collapses (the moral hazard threshold).
\item \textbf{Asymmetric mutual contract veto.}\ Inter-civ insurance
contracts (Z) are rejected if projected to be persistently
asymmetric (one party always contributing, never drawing). This
prevents predatory collective structures.
\item \textbf{Drift threshold.}\ No configuration is admissible if
the projected rescue\_dependency trajectory (AB) exceeds
\texttt{rescue\_dependency\_ceiling}. This prevents the
collective from absorbing all individual agency.
\end{enumerate}

These four veto rules instantiate the NRMO architectural principle
at the collective scale: the search engine operates only within the
admissible set.

\subsection{Honest design constraints}

\paragraph{Computational cost.}
v7.1's Collective Engine performs MC rollout each generation for
each civilisation. Cost scales as $C_{\text{candidates}} \times
\text{rollout\_depth} \times \text{rollout\_repeats} \times
N_{\text{civilisations}} \times N_{\text{generations}}$. With
candidate count $\approx 12$, depth $6$, repeats $4$, this is
$\approx 1{,}728$ collective transitions per civ per gen, on top
of individual $\Omega$ Full. Small scale runs are projected to take
50--80\% longer than v7.0.

\paragraph{Score function design.}
The collective rollout score combines continuation rate (avoid civ
collapse) and Shannon strategy diversity (avoid monoculture):
\[
\mathit{score}_{\text{collective}} =
\alpha \cdot \mathrm{continuation\_rate}_{\text{civ}} +
\beta \cdot H_{\text{strategy}}(p) +
\gamma \cdot \mathrm{cohesion}
\]
with weights $(\alpha, \beta, \gamma)$ that may differ by world
archetype. The Shannon entropy term operationalises the user
requirement ``lineage diversity must be preserved as an objective,
not only avoided as a failure''.

\paragraph{No simplification.}
v7.1 is implemented with all components in place: predictive
trigger, full triage with learned weights, dynamic reallocation,
inter-civ contract negotiation, candidate exploration, drift
control, four veto rules. Per the user directive, ``simple
implementation prohibited.''
